% Sitography: every link in one place, clickable.
% Verified July 2026. Links rot; the repository issue tracker is the place to
% report a dead one.

%--------------------------------------------------------------------------
\section{How to Use This Chapter} \label{sec:sito-how}

Every project, tool, specification and paper referenced in this book is
collected here with a clickable link. If you are reading the PDF, the titles
are live: click them.

The list is deliberately opinionated. Where two resources cover the same
ground, the more useful one is listed and the other is not. Where a project is
promising but not yet usable, that is stated.

\begin{remark}
Links rot. This sitography was verified in July 2026. The canonical, updatable
version lives in the repository, and a dead link is a legitimate issue to
open:
\href{https://github.com/elRaulito/ZK-from-zero-on-Cardano}{github.com/elRaulito/ZK-from-zero-on-Cardano}
\end{remark}

%--------------------------------------------------------------------------
\section{This Book} \label{sec:sito-book}

\begin{description}
  \item[\href{https://github.com/elRaulito/ZK-from-zero-on-Cardano}{ZK from Zero on Cardano}]
    The repository. LaTeX source, the \texttt{zk-password} reference
    implementation, the benchmarks behind Chapter~\ref{ch:tornado}, and the
    rendered PDFs on the \texttt{rendered} branch.

  \item[\href{https://projectcatalyst.io/funds/14/cardano-open-ecosystem/zk-from-zero-on-cardano-ebook-by-elraulito}{Catalyst proposal (Fund 14, \#1400129)}]
    The original proposal, milestone reports, and close-out documentation.

  \item[\href{https://www.raul.it/}{raul.it}]
    The author.
\end{description}

%--------------------------------------------------------------------------
\section{Cardano ZK Projects} \label{sec:sito-cardano-zk}

This is the complete set of open-source ZK work on Cardano known to the author
at the time of writing. It is short enough to read in a week. That is either
intimidating or an opportunity.

\subsection*{Applications}

\begin{description}
  \item[\href{https://github.com/leobel/janus-wallet}{Janus Wallet}]
    The most complete open-source ZK application on Cardano. A wallet where
    the user authenticates with a password rather than a seed phrase, proving
    knowledge of it via a Groth16 circuit. Circom, Poseidon, BLS12-381, Aiken
    validators, and a Node.js proving service. \textbf{Read the circuit and
    the validators together}; it is the best available worked example of how a
    real ZK product is assembled. Discussed in
    Chapters~\ref{ch:circuits}, \ref{ch:password} and~\ref{ch:offchain}.

  \item[\href{https://github.com/Modulo-P/Cardano-Semaphore}{Cardano-Semaphore}]
    A port of Semaphore, the anonymous-signalling and group-membership
    protocol, by Modulo-P. Semaphore has the same Merkle-tree-plus-nullifier
    shape as the mixer analysed in Chapter~\ref{ch:tornado}, which makes this
    the project to watch for progress on the on-chain accumulator problem.

  \item[\href{https://midnight.network/}{Midnight}]
    Cardano's ZK partner chain. Dual-ledger architecture separating public
    from private data, with selective disclosure as a first-class feature.
    Federated mainnet launched March 2026. Cardano's institutional answer to
    the applications Chapter~\ref{ch:tornado} shows the L1 cannot host.
\end{description}

\subsection*{Libraries and Tooling}

\begin{description}
  \item[\href{https://github.com/eryxcoop/cardano-zk-aiken}{aiken-zk (Eryx)}]
    A language extension and CLI that generates the Circom circuit, the Aiken
    verifier, and even the MeshJS redeemer from an \texttt{offchain}
    annotation in your validator. Catalyst Fund 13. The most promising
    developer-experience work in the space. Note the honest caveat on its
    Merkle tree hash, discussed in Section~\ref{sec:ch7-aikenzk}.

  \item[\href{https://github.com/Modulo-P/ak-381}{ak-381 (Modulo-P)}]
    A reusable Aiken implementation of a Groth16 verifier over BLS12-381. If
    you would rather import a verifier than write the one in
    Chapter~\ref{ch:password}, start here.

  \item[\href{https://github.com/cardano-foundation/bls}{cardano-foundation/bls}]
    A "gym" of BLS12-381 examples in Aiken, from the Cardano Foundation. The
    place to look for primitive-level patterns.

  \item[\href{https://github.com/jmagan/poseidon-bls12381}{poseidon-bls12381}]
    A TypeScript Poseidon implementation targeting BLS12-381. Useful
    off-chain, and a solution to the \texttt{circomlibjs} field-mismatch trap
    described in Section~\ref{sec:ch7-provers}.
\end{description}

\subsection*{Background Reading}

\begin{description}
  \item[\href{https://cardanofoundation.org/blog/aiken-primitives-explained}{Aiken's BLS12-381 Primitives Explained}]
    Cardano Foundation walkthrough of what the pairing builtins make possible.

  \item[\href{https://iohk.io/en/blog/posts/2024/02/12/unlocking-more-opportunities-with-plutus-v3/}{Unlocking More Opportunities with Plutus V3}]
    The IOG post introducing the primitives this entire book depends on.
\end{description}

%--------------------------------------------------------------------------
\section{Circuits and Proving} \label{sec:sito-circuits}

\begin{description}
  \item[\href{https://docs.circom.io/}{Circom documentation}]
    The circuit language. Start with \emph{Getting Started}, then the
    \emph{Circom Language} reference. Remember
    \texttt{-{}-prime bls12381}.

  \item[\href{https://github.com/iden3/circomlib}{circomlib}]
    The primitive circuit library: Poseidon, comparators, bit decomposition,
    Merkle proofs. Read \texttt{poseidon.circom} against the round counts in
    Section~\ref{sec:ch8-wall1}.

  \item[\href{https://github.com/iden3/snarkjs}{snarkjs}]
    Trusted setup, proving, and verification. The reference implementation and
    the one this book uses throughout.

  \item[\href{https://github.com/iden3/rapidsnark}{rapidsnark}]
    A C++ prover, 5 to 10 times faster than snarkjs, reading the same
    \texttt{.zkey} files. Check its BLS12-381 support before committing.

  \item[\href{https://github.com/arkworks-rs/circom-compat}{ark-circom}]
    Drive Circom circuits from Rust via arkworks. First-class BLS12-381
    support.

  \item[\href{https://github.com/Consensys/gnark}{gnark}]
    A Go proving library with its own circuit DSL. No Cardano verifier, but
    good for understanding proving systems from the inside.

  \item[\href{https://noir-lang.org/}{Noir}]
    A far nicer ZK language than Circom. Listed so you know it exists and know
    why this book does not use it: its backends have no Aiken verifier.
\end{description}

%--------------------------------------------------------------------------
\section{Cardano Development} \label{sec:sito-cardano-dev}

\begin{description}
  \item[\href{https://aiken-lang.org/}{Aiken}]
    The on-chain language. The
    \href{https://aiken-lang.github.io/stdlib/}{standard library} docs are
    where the \texttt{bls12\_381} modules live, and reading the
    \texttt{scalar} source explains a great deal about why
    Chapter~\ref{ch:tornado} reaches the conclusion it does.

  \item[\href{https://meshjs.dev/}{MeshJS}]
    The TypeScript SDK used for the transactions in
    Chapter~\ref{ch:offchain}. See
    \href{https://meshjs.dev/apis/txbuilder/smart-contracts}{Smart Contract Interactions}
    for the \texttt{MeshTxBuilder} patterns.

  \item[\href{https://github.com/no-witness-labs/lucid-evolution}{Lucid Evolution / Evolution SDK}]
    The main alternative to MeshJS. Same job, different ergonomics. Originally
    from Anastasia Labs, now maintained by No Witness Labs and being renamed to
    Evolution SDK, so expect the package names to move.

  \item[\href{https://blockfrost.io/}{Blockfrost}]
    Hosted chain API. Free tier is sufficient for everything in this book.

  \item[\href{https://www.koios.rest/}{Koios}]
    A public, no-key-required Cardano API. The live protocol parameters quoted
    in Chapters~\ref{ch:offchain} and~\ref{ch:tornado} came from here.

  \item[\href{https://gomaestro.org/}{Maestro}]
    Another provider, with indexing features useful once you outgrow simple
    UTxO queries.

  \item[\href{https://cardanoscan.io/}{Cardanoscan}] and
    \href{https://cexplorer.io/}{cexplorer}\textbf{.}
    Explorers. Useful for inspecting datums and script execution units on a
    submitted transaction.

  \item[\href{https://docs.cardano.org/about-cardano/explore-more/parameter-guide}{Protocol parameters guide}]
    What each parameter means. For live \emph{values}, query Koios or a node.
\end{description}

%--------------------------------------------------------------------------
\section{Cardano Improvement Proposals} \label{sec:sito-cips}

How cryptographic primitives actually arrive on Cardano. Chapter
\ref{ch:future} argues that a Poseidon builtin belongs on this list and is not
yet on it.

\begin{description}
  \item[\href{https://cips.cardano.org/cip/CIP-0381}{CIP-381}]
    The BLS12-381 primitives themselves. The foundation of every ZK contract
    on Cardano today.

  \item[\href{https://cips.cardano.org/cip/CIP-0101}{CIP-101}]
    Keccak-256, which enables Ethereum signature verification in scripts.

  \item[\href{https://cips.cardano.org/cip/CIP-0109}{CIP-109}]
    \texttt{modularExponentiation}, and with it on-chain modular inverses.
    Still at \emph{Proposed} status. The most directly relevant open proposal
    to the cost floor measured in Section~\ref{sec:ch8-wall1}.

  \item[\href{https://cips.cardano.org/cip/CIP-0127}{CIP-127}]
    RIPEMD-160.

  \item[\href{https://github.com/cardano-foundation/CIPs}{The CIP repository}]
    Where to propose one, and where to read the discussion on proposals that
    have not landed.
\end{description}

%--------------------------------------------------------------------------
\section{Papers and Foundations} \label{sec:sito-papers}

The two worth reading in full rather than in summary are marked.

\begin{description}
  \item[\href{https://dl.acm.org/doi/10.1145/22145.22178}{Goldwasser, Micali, Rackoff (1985)}]
    \emph{The Knowledge Complexity of Interactive Proof Systems.} The paper
    that defines what a zero-knowledge proof \emph{is}.
    \textbf{Read this one in full.}

  \item[\href{https://eprint.iacr.org/2016/260}{Groth (2016)}]
    \emph{On the Size of Pairing-Based Non-interactive Arguments.} The
    construction whose verification equation we implemented in
    Chapter~\ref{ch:password}. \textbf{Read this one too.}

  \item[\href{https://eprint.iacr.org/2019/458}{Grassi et al.\ (2019)}]
    \emph{Poseidon: A New Hash Function for Zero-Knowledge Proof Systems.}
    The round structure costed in Section~\ref{sec:ch8-wall1}.

  \item[\href{https://eprint.iacr.org/2019/953}{Gabizon, Williamson, Ciobotaru (2019)}]
    \emph{PLONK.} The universal-setup construction that Chapter
    \ref{ch:future} argues Cardano should eventually support.

  \item[\href{https://eprint.iacr.org/2018/046}{Ben-Sasson et al.\ (2018)}]
    \emph{Scalable, Transparent, and Post-quantum Secure Computational
    Integrity.} STARKs. Relevant to the speculative research direction in
    Section~\ref{sec:ch9-missing}.

  \item[\href{https://zkproof.org/}{ZKProof}]
    Standards work, plus a curated reading path into the wider literature.

  \item[\href{https://zcash.github.io/halo2/}{The Halo2 Book}]
    For understanding what a no-trusted-setup proving system requires of its
    verifier.
\end{description}

%--------------------------------------------------------------------------
\section{Reference Implementations Worth Reading} \label{sec:sito-reference}

Not Cardano projects, but the canonical implementations of the ideas this book
discusses. Reading them is how you learn what a mature ZK codebase looks like.

\begin{description}
  \item[\href{https://github.com/tornadocash/tornado-core}{tornado-core}]
    The original mixer contracts and circuits reverse-engineered in
    Chapter~\ref{ch:tornado}. Listed for study. As that chapter sets out at
    length, the protocol was sanctioned in 2022 and this book does not
    recommend deploying anything derived from it.

  \item[\href{https://github.com/semaphore-protocol/semaphore}{Semaphore}]
    Anonymous signalling on Ethereum, and the upstream of Cardano-Semaphore.
    A cleaner introduction to Merkle-tree-plus-nullifier design than Tornado.

  \item[\href{https://github.com/zcash/zcash}{Zcash}]
    The first large-scale zk-SNARK deployment, discussed in
    Chapter~\ref{ch:web3}.

  \item[\href{https://github.com/MinaProtocol/mina}{Mina}]
    Recursive proof composition taken to its logical conclusion: a whole
    blockchain compressed into one constant-size proof.
\end{description}

%--------------------------------------------------------------------------
\section{Where to Ask} \label{sec:sito-community}

\begin{description}
  \item[\href{https://discord.gg/aiken}{Aiken Discord}]
    The most useful place for Aiken and on-chain questions, including the
    BLS12-381 builtins.

  \item[\href{https://forum.cardano.org/}{Cardano Forum}]
    Longer-form technical discussion. The \emph{Smart Contracts} category is
    where execution-unit questions get answered.

  \item[\href{https://cardano.stackexchange.com/}{Cardano Stack Exchange}]
    Best for specific, answerable questions.

  \item[\href{https://t.me/CardanoDevelopersOfficial}{Cardano Developers (Telegram)}]
    Fast, informal, good for unblocking.

  \item[\href{https://zkhack.dev/}{ZK Hack}]
    Not Cardano-specific, and the best community for questions about the
    cryptography itself: study groups, puzzle challenges, and talks. The
    \href{https://zeroknowledge.fm/}{Zero Knowledge podcast} is the single
    best way to keep up with the field.
\end{description}

\begin{highlight}
If you build something with this book, the most useful thing you can do is
write down what broke. The failure tables in Chapter~\ref{ch:offchain} and the
measurements in Chapter~\ref{ch:tornado} exist because the problems were
documented rather than merely worked around. This sitography is short. Help
make it longer.
\end{highlight}
